\documentclass[12pt]{article}
\usepackage{enumitem}
\usepackage[top=0.6in,bottom=0.6in,right=0.7in,left=0.7in]{geometry}
\usepackage[LGR, T1]{fontenc}
\usepackage[utf8]{inputenc}
\usepackage[greek, english]{babel}
\usepackage{alphabeta}
\usepackage{graphicx}
\usepackage{subfig}
\usepackage{listings}
\usepackage[breaklinks]{hyperref}
\usepackage[dvipsnames]{xcolor}
\usepackage{amsmath}
\usepackage{listings}
\usepackage{ragged2e}
\usepackage{titlepic}
\usepackage{hyperref}
\usepackage{booktabs}

\usepackage[giveninits=true, style=ieee]{biblatex}
\addbibresource{references.bib}
\renewbibmacro{in:}{}

\graphicspath{ {images/} }
%%%%%%%%%%%%%%%%%%%%%%%%%%%%%%%%%%%%%%%%%

\begin{document}
\begin{titlepage}
% \pagecolor{blue!10}
\begin{center}

\textsc{\Large \textcolor{black}{\textbf{ΕΘΝΙΚΟ ΜΕΤΣΟΒΙΟ ΠΟΛΥΤΕΧΝΕΙΟ}}}\\[0.4cm]	

{\huge {Σχολή Ηλεκτρολόγων Μηχανικών και Μηχανικών Υπολογιστών} \\[1cm] }
{\huge {School of Electrical and Computer Engineering} \\[1cm] }


\includegraphics[scale=0.3]{NTUA_Logo.png}\\[1cm]

\textsc{\Large \textbf{Μη-Γραμμική Δυναμική Συστημάτων και Ταλαντώσεις}}\\[0.4cm]
\textsc{\Large \textbf{Non-Linear Dynamics and Oscillations}}\\[0.4cm]


\rule{\linewidth}{0.7mm} \\[0.4cm]
{ \huge \bfseries\color{black!70!black} Bifurcation Analysis of Spiking Nonlinear Opinion Dynamics}
\rule{\linewidth}{0.7mm} \\[2cm]

\begin{tabular}{l}

\Large Σπυρίδων Γιάνι (03120026)\\[0.4cm]

\end{tabular}

\vfill


{\large \color{black!80!black}{\textbf{Athens, Greece}}\\[0.2cm] \color{blue!90!black}23 February 2026}

\end{center}
\end{titlepage}

\newpage
\tableofcontents

\newpage

\section{Overview of the Nonlinear Opinion Dynamics Model}

\subsubsection*{Introduction}
\label{intro}

Nonlinear Opinion Dynamics (NOD) models are used to describe how opinions evolve over time within a population. They are fast and flexible, 
in the sense that on one hand they diverge quickly from indecision, while on the other hand their sensitivity is tunable. 
The equation defining it \cite{Cathcart2024} is the following:
\begin{equation}
    \dot{z} = -d z + tanh(u (a z) + b)
    \label{nod}
\end{equation}
where: 
\begin{itemize}
    \item $z(t)$ is the system's opinion, about 2 options (positive corresponds to option 1, negative to option 2). The magnitude $\lvert z \rvert$ 
    shows the magnitude of the agent's opinion on the repsective option
    \item $u(t)$ is the agent's sensitivity to its observations and it is implemented as a feedback gain, where $u = u_0 + K_u z^2$
    \item $b(t)$ is the external input
    \item $d$ is the damping coefficient of $z(t)$
    \item $a$ is the amplification coefficient of $z(t)$
    \item $K_u$ is the feedback gain
\end{itemize}

In \cite{Cathcart2024}, the dynamics of the NOD system are studied, showing that the gain $K_u$ plays a fundamental role: (Proposition 1)
large values of $K_u$ increase the strength of non-neutral ($z \neq 0$) opinions. Specifically, at a critical value $u_0 = u_0^{*}$, there is a: 
\begin{itemize}
    \item supercritical pitchfork bifurcation, for $K_u < \frac{d^3}{3a}$
    \item subcritical pitchfork bifurcation, for $K_u > \frac{d^3}{3a}$
\end{itemize}

In the case of the subcritical pitchfork, for $u_0 > u_0^{*}$ the system has three equilibria: the origin and two stable non-neutral, which are 
also symmetric with respect to $z = 0$. These two stable equilibria are stronger for large values of $K_u$ and allow the model 
to converge in a quick and stable manner to a specific option. 

\subsubsection*{Limitations and the S-NOD Model}

The NOD model described by \ref{nod} has a serious limitation: for large values of $K_u$, the system's opinion $z(t)$
is not affected, unless very large inputs $b(t)$ are applied. This is an undesirable property in rapidly changing systems, where 
agents must be able to quickly adapt to new information. Simply reducing $K_u$ would not be a preferable solution, as the decision 
would not be as fast and robust. This trade-off is resolved with the Spiking-Nonlinear Opinion Dynamics (S-NOD) model, an alternate version to 
\ref{nod}, which offers more flexibility. 

It is defined by the following set of equations: 
\begin{align}
    \dot{z} & = -d z + tanh((k z^2 + \mu_0 - s)\alpha z + b) \\
    \dot{s} & = \epsilon (k_s z^4 - s)
    \label{snod}
\end{align}
where the new variables and parameters introduced are: 
\begin{itemize}
    \item $s(t)$ is the agent's slow recovery state, which provides negative feedback to the opinion variable 
    \item $k$ is the sensitivity feedback gain
    \item $m_0$ is the basal leve of the sensitivity 
    \item $\epsilon$ is the timescale separation between the dynamics of $z(t)$ and $s(t)$
    \item $k_s$ is the gain of the negative feedback 
\end{itemize}

The S-NOD system maintains the advantages of the NOD model, in the sense that while $z(t)$ increases, the positive feedback causes 
the model to deviate quickly from the equilibrium and settle on a specific opinion. However, the slow negative feedback, introduced 
by the state $s(t)$, causes the sensitivity of the agent to decrease over time, allowing the model to return to a neutral state and be more 
receptive to new inputs. 

Studying the dynamics of the system, this "back-and-forth" behaviour of the opinion variable $z(t)$ is evident through the 
formation of limit cycles in the phase plane. The goal of this project is to prove the emergence of these limit cycles and analyse 
the bifurcations causing them. 

\section{Phase Space Analysis}

In this section, we will examine how varying a single system parameter (while keeping the other parameters fixed), 
by studying trajectories in the phase plane. To this end, a Python code has been developed, which, given the
2-dimensional dynamical system, plots the phase portrait of the system in the (z, s) plane.
\footnote{The Python code, titled phase\_space.py, is available in this Github repository: 
\url{https://github.com/SpyrosJani/snod_project}, along with the report in PDF format and its LaTeX source code.} 

\subsection{Case A: Input-Driven Spiking}

The first case is concerned with how the constant external input $b$ affectes the S-NOD system \ref{snod}. Following 
\cite{Belaustegui2025}, the other parameters are fixed in the following values: 
\begin{center}
    \begin{math}
        \alpha = 1, d = 1, k = 3, k_s = 22, \epsilon = 0.05, \mu_0 = 0.8 
    \end{math}
\end{center}
Here we will consider only values $b \ge 0$, as the case for $b < 0$ is identical, 
due to the odd symmetry of the model.  

In dynamical systems, it is common practice to begin the analysis by locating equilibria. The analytical expressions 
of the equilibria are not easily computable, but Proposition 3.2 of \cite{Belaustegui2025} proves that the S-NOD 
system has a unique equilibrium for any value of $b$. Furthermore, Theorem 3.7 of \cite{Belaustegui2025} describes 
the bifurcations happening in the model: 
\begin{center}
    \textit{For certain conditions between the rest of the parameters (which can be easily verified that hold true for the 
    chosen values), the unique equilibrium undergoes a Hopf bifurcation at two critical input values, $b^{*}$ and $b^{\dag}$, 
    where $0 < b^{*} < b^{\dag}$}
\end{center}
Of course, due to the odd symmetry, for negative values of the input, the two critical values are $-b^{*}$ and $-b^{\dag}$.

Combined with Theorem 3.8, for an input below the \emph{spiking threshold} $b^{*}$ the system has only 
an equilibrium, while above $b^{*}$ limit cycles are created, which correspond to spiking of the opinion variable $z(t)$.
Figure 2 in \cite{Belaustegui2025} suggests that the spiking threshold occurs at $b = 0.1$, therefore, in this analysis 
we will be input values above and below this threshold. 

\subsubsection*{Below the Threshold}

For this case, we choose the value $b = 0.02$. It is often helpful to plot the equilibria 
and the nullclines, before any trajectories. This is shown in figure \ref{below_nullclines}, 
where the equilibrium is the intersection of the 2 nullclines.
In figure \ref{below}, four trajectories of the S-NOD system have been plotted, 
with randomly chosen initial conditions. All trajectories approach the unique equilibrium 
asymptotically, along the z-nullcline, which suggests that the equilibrium is stable.
\begin{figure}
    \centering    
    \includegraphics[width = 0.6\textwidth]{case_a_phase_plane_below_without.png}
    \caption{The s-nullcline (light grey) has the shape of a parabola, while the z-nullcline (dark grey) 
    comprises of 2 branches, both of which asymptotically approach $z=0$ and $z = \pm 1$}
    \label{below_nullclines}
\end{figure}

\begin{figure}
    \centering    
    \includegraphics[width = 0.6\textwidth]{case_a_phase_plane_below.png}
    \caption{The phase portrait of the S-NOD system for $b = 0.02 < b^{*}$. All trajectories converge to the unique equilibrium, 
    which appears to be stable.}
    \label{below}
\end{figure}

\subsubsection*{Above the Threshold}

We examine how the system behaves for $b = 0.2 > b^{*}$. Again, we plot the nullclines and the equilibrium
in figure \ref{above_nullclines}. With four randomly chosen initial conditions, 
we plot the trajectories in figure \ref{above}. It is clear that the trajectories converge to a limit cycle, 
following the z-nullcline. Thus, \emph{opinion spiking} occurs, i.e. the opinion variable is oscillating with 
respect to time.
The limit cycle appears to be stable and it encircles the unique equilibrium. 
\footnote{If the input $b(t)$ becomes smaller than $-b^{*}$, then a limit cycle is also formed, but it will correspond to negative values 
of the opinion variable $z(t)$, i.e. in option 2.}
\begin{figure}
    \centering    
    \includegraphics[width = 0.6\textwidth]{case_a_phase_plane_above_without.png}
    \caption{The s-nullcline (light grey) and the z-nullcline (dark grey) of the 
    system, as before.}
    \label{above_nullclines}
\end{figure}

\begin{figure}
    \centering    
    \includegraphics[width = 0.6\textwidth]{case_a_phase_plane_above.png}
    \caption{The phase portrait of the S-NOD system for $b = 0.2 < b^{*}$. All trajectories converge 
    to a limit cycle, which is clearly identified.}
    \label{above}
\end{figure}

\subsubsection*{Interpretation}

The results above show us that when the external input is not strong enough, the system (in the "long run") remains deadlocked (in other 
words indecisive). However, when the input is stronger, the system converges to a specific option. This is the same behaviour as the NOD 
system, studied in section \ref{intro}. The crucial difference is that now, the robustness of the decided comes in an excitable form, i.e. 
through limit cycles. This allows the opinion variable $z(t)$ to oscillate and be flexible enough to change, without requiring the external 
input to become extremely large.

\subsection{Case B: Deadlock Breaking via Basal Sensitivity}

Now the input $b$ is kept constant, but the basal sensitivity parameter $\mu_0$ is varied. Specifically, the parameters 
are kept fixed as such: 
\begin{center}
    \begin{math}
        \alpha = 1, d = 1, k = 2.3, k_s = 16, \epsilon = 0.1, b = 0
    \end{math}
\end{center} 
following \cite{Belaustegui2025}. 

According to \cite{Cathcart2024}, because $b = 0$, the origin $(0, 0)$ is always an equilibrium. For a basal sensitivity below the 
critical value $\mu_0^{*}$, the origin is stable, while at $\mu_0 = \mu_0^{*}$ a subcritical pitchfork bifurcation occurs, 
causing it to lose stability \cite{Belaustegui2025}. This critical value is: 
\begin{center}
    \begin{math}
        \mu_0^{*} = \frac{d}{\alpha} = 1
    \end{math}
\end{center} 
and according to Proposition 4.2 and Theorem 4.3 of \cite{Belaustegui2025}, for $\mu_0 > \mu_0^{*}$ 
two more equilibria are created, symmetric around $z = 0$, while two symmetric limit cycles also appear, 
encircling these two non-neutral equilibria. We will verify these results by plotting the phase portraits of the system. 

\subsubsection*{Below the Critical Basal Sensitivity}

We choose a basal sensitivity value of $\mu_0 = 0.8 < \mu_0^{*}$. The nullclines and equilibrium are shown in figure \ref{below_b}.
The s-nullcline is, as before, a parabola. However, the z-nullcline contains two branches: for $z \neq 0$ it is a sigmoid-like curve, bounded 
between -1 and 1. For $z = 0$, it is clear from the definition of the system that any value of $s$ satisfies the z-nullcline equation ($\dot{z} = 0$)
So, the second branch of the nullcline will be the horizontal line $z = 0$, which is the only part of the z-nullcline that intersects 
with the s-nullcline, as it is clearly seen in figure \ref{below_b}.

As before, we choose four random initial conditions and plot the trajectories in figure \ref{below_b}. They all converge 
to the origin, following the $z = 0$ branch of the z-nullcline asymptotically, causing the system to be deadlocked and to not make 
any decisions.
\begin{figure}
    \centering    
    \includegraphics[width = 0.6\textwidth]{case_b_phase_plane_below.png}
    \caption{The phase portrait of the S-NOD system for $\mu_0 = 0.8$. The z-nullcline (dark grey) 
    and the s-nullcline (light grey) intersect only at the origin, which appears to be a stable equilibrium.}
    \label{below_b}
\end{figure}

\subsubsection*{Above the Critical Basal Sensitivity}

In this occasion, we choose $$\mu_0 = 1.2 > \mu_0^{*}$$. The nullclines are shown in figure \ref{above_nullclines_b} and it is obvious that two 
new equilibria have formed, symmetric around the horizontal axis. 
\begin{figure}
    \centering    
    \includegraphics[width = 0.6\textwidth]{case_b_phase_plane_above_without.png}
    \caption{The nullclines and the equilibria of the S-NOD system for basal sensitivity above the critical. The z-nullcline (dark grey) 
    and the s-nullcline (light grey) intersect not only at the origin, but at two more points, which are symmetric.}
    \label{above_nullclines_b}
\end{figure}

In figure \ref{above_b}, we plot the phase portrait of the system. As we can see, the deadlock breaks and trajectories do not converge to the 
neutral solution. Specifically, any initially positive opinion (no matter how small) will converge to a stable limit cycle, which is 
strictly positive (for values of $z(t)$) and encircles the positive unstable equilibrium. Respectively, a negative opinion always ends up 
in a limit cycle, which is strictly negative (for values of $z(t)$). Furthermore, the two limit cycles are symmetrical. 

\begin{figure}
    \centering
    \begin{minipage}{0.49\textwidth}
        \centering
        \includegraphics[width = \textwidth]{case_b_phase_plane_above_v1.png}
    \end{minipage}
    \begin{minipage}{0.49\textwidth}
        \centering
        \includegraphics[width = \textwidth]{case_b_phase_plane_above_v2.png}
    \end{minipage}
    \caption{The phase portrait of the S-NOD system for $\mu_0 = 1.2$. In the left figure, the initial conditions of the trajectories 
    have been chosen to be outside the limit cycles, while in the right figure, they have been chosen to be inside. In both cases, 
    the trajectories converge to the same limit cycles, depending on whether the initial opinion variable $z(0)$ was positive or negative.}
    \label{above_b}
\end{figure}

\subsubsection*{Interpretation}

The analysis above shows that if the system is not sensitive enough, if no external input is applied, it remains neutral, an expected behaviour. 
If the sensitivity grows, with no external input, even small perturbations (positive or negative), push the model to make a non-neutral 
decision (towards option 1 or option 2, respectively). Nonetheless, the excitability offered by the S-NOD model (again via the formation 
of limit cycles) allow the agent to be flexible enough to respond effectively to external changes in the environment. 

\newpage
\section{Bifurcation Analysis}

In this section, we will numerically verify the bifurcations described in \cite{Belaustegui2025}. This will be done using the MatCont Toolbox
\footnote{\url{https://matcont.sourceforge.io/}}, 
a MatLab-based software for numerical continuation and bifurcation study of dynamical systems. We will examine the same two cases as before,
by calculating equilibria continuation, limit cycles and the resulting opinion spikes.

\subsection{Case A: Input-Driven Spiking}

In MatCont, we define the S-NOD system \ref{snod} and set the parameters as in the previous section. We perform an equilibrium continuation
with respect to the input $b$, starting from $b = 0$. The initial condition is set to be the origin, since it is the only equilibrium 
for $b = 0$. After one forward and one backward continuation, we obtain the bifurcation diagram shown in figure \ref{bif_a}.
\begin{figure}[ht]
    \centering    
    \includegraphics[width = 0.6\textwidth]{case_a_equilibria_cont.png}
    \caption{Bifurcation diagram of the S-NOD system with respect to the input $b$. The blue line indicates that the equilibrium is 
    stable and the purple line that it is unstable. Hopf bifurcations are marked with the letter "H".}
    \label{bif_a}
\end{figure}

As predicted in Theorem 3.7 of \cite{Belaustegui2025}, four Hopf bifurcations occur in total, in symmetrical pairs. 
Two of them occur for positive values of $b$ (at $b^{*}$ and $b^{\dag}$), and two at their negative counterparts. 
For input values between $b^{*}$ and $b^{\dag}$, the equilibrium is unstable and as we have seen before, trajectories converge 
to stable limit cycles. It is interesting to note, that this bifurcation diagram suggests that very large values of the input 
eliminate opinion spikes and cause the system to converge strongly and robustly to a specific option, without oscillations. 

We can also verify that Hopf bifurcations occur by observing the eigenvalues of the Jacobian matrix of the system at the 
bifurcations, as seen in table \ref{hopf_eigenvalues}. In both cases, at the bifurcation points, 
the eigenvalues are purely imaginary (with zero real part), which is the defining property of Hopf bifurcations.
\begin{table}[ht]
    \centering
    %\begin{tabular}{| p{0.1\linewidth} | p{0.3\linewidth} | p{0.3\linewidth} |}
    \begin{tabular}{c|c|c}
        \toprule
        \textbf{$b$} & \textbf{Re($\lambda$)} & \textbf{Im($\lambda$)} \\
        \midrule
        $0.0250452$ & $6.1826 \cdot 10^{-7}$ & $0.051989$\\
        %\hline 
         & $6.1826 \cdot 10^{-7}$ & $-0.051989$\\
        \hline 
        $0.7909516$ & $-1.0309 \cdot 10^{-6}$ & $0.528419$\\
        %\hline 
         & $-1.0309 \cdot 10^{-6}$ & $-0.528419$\\
        %\hline 
        \bottomrule
    \end{tabular}
    \caption{Critical values of the input $b$ at which Hopf bifurcations occur, 
    along with the eigenvalues of the Jacobian matrix of the system}
    \label{hopf_eigenvalues}
\end{table}

Finally, we can perform a limit cycle continuation, starting from the Hopf bifurcation at $b^{*}$, 
which results in the diagram shown in figure \ref{limit_cycle_cont}. From both diagrams it interesting to note 
that while we approach the second Hopf bifurcation at $b^{\dag}$, the period of the limit cycle decreases significantly, 
while the magnitude of the opinion variable $z(t)$ exhibits significantly smaller spiking. 
This suggests that the system is less excitable and more robust to changes in the input, as we approach $b^{\dag}$.
Note that these diagrams have occured after the 'MaxStepSize' parameter of MatCont has bee increased to 2, 
so that the appearing limit cycles are distinctly visible.
\begin{figure}
    \centering
    \begin{minipage}{0.49\textwidth}
        \centering
        \includegraphics[width = \textwidth]{case_b_limit_cycle_cont_a.png}
    \end{minipage}
    \begin{minipage}{0.49\textwidth}
        \centering
        \includegraphics[width = \textwidth]{case_b_limit_cycle_cont_b.png}
    \end{minipage}
    \caption{Result of limit cycle continuation of the system, starting from the Hopf bifurcation at $b^{*}$. 
    The left figure shows the opinion variable $z(t)$ with respect to the input $b$, 
    while the right figure shows the period of the limit cycle with respect to $b$.}
    \label{limit_cycle_cont}
\end{figure}

\subsection{Case B: Deadlock Breaking via Basal Sensitivity}

Fixing the parameters as before and choosing again the inital point to be the origin, we perform an equilibrium continuation 
with respect to the basal sensitivity $\mu_0$, starting from $\mu_0 = 0.7$.
After one forward and one backward continuation, a "Branch Point" bifurcation is detected at $\mu_0^{*} = 1$, 
which is to be expected (a subcritical pitchfork bifurcation is supposed to occur). 
We choose this point as a new initial condition and perform a second equilibrium continuation,
which results in the bifurcation diagram shown in figure \ref{bif_b}.
\begin{figure}[ht]
    \centering    
    \includegraphics[width = 0.6\textwidth]{case_b_equilibria_cont.png}
    \caption{Bifurcation diagram of the S-NOD system with respect to the basal sensitivity $\mu_0$. The blue line indicates that the equilibrium is 
    stable and the purple line that it is unstable. The "BP" marker indicates the Branch Point bifurcation at $\mu_0^{*} = 1$.
    The "H" here corresponds to a "Neutral Saddle" bifurcation, not to a Hopf. The "LP" marker 
    indicates the "Limit Point", which is a saddle-node bifurcation.}
    \label{bif_b}
\end{figure}

The diagram confirms the theoretical predictions and the phase portraits. For small basal sensitivity, the only 
equilibrium is the origin, which is stable. At $\mu_0^{*} = 1$, a subcritical pitchfork bifurcation occurs, 
making the origin unstable. For $\mu_0 > \mu_0^{*}$, two new non-zero, unstable equilibria are created, symmetric around $z = 0$, 
through a saddle-node bifurcation. Each of these equilibria is encircled by a stable limit cycle, which corresponds to opinion spiking, 
as we have seen in the phase space analysis.






\newpage
\printbibliography[heading=bibintoc]

\end{document}